\documentclass[11pt,aspectratio=169]{ctexbeamer}

\usetheme{default}
\setbeamertemplate{navigation symbols}{}

\usepackage{amsmath}
\usepackage{booktabs}

% Projector-safe palette: high contrast on common lecture-room projectors.
\definecolor{DeepBlue}{RGB}{11,58,102}
\definecolor{SoftInk}{RGB}{31,41,51}
\definecolor{PaperBg}{RGB}{255,255,255}
\definecolor{PanelBlue}{RGB}{226,237,247}

\setbeamercolor{normal text}{fg=SoftInk,bg=PaperBg}
\setbeamercolor{structure}{fg=DeepBlue}
\setbeamercolor{frametitle}{fg=DeepBlue,bg=PaperBg}
\setbeamercolor{title}{fg=DeepBlue}
\setbeamercolor{block title}{fg=DeepBlue,bg=PanelBlue}
\setbeamercolor{block body}{fg=SoftInk,bg=white}

\newcommand{\Core}[1]{
  \begin{block}{核心结论}
    #1
  \end{block}
}

\newcommand{\LectureMeta}[2]{
  \title{大模型推理系统与实践}
  \subtitle{#1}
  \author{张书豪}
  \institute{华中科技大学计算机学院}
  \date{2027 · #2}
}

\newcommand{\PaperRef}[3]{
  \item \textbf{#1}（#2）\hfill\href{#3}{原文链接}
}

\newcommand{\ResearchSources}[1]{
  \begin{frame}{本讲参考资料}
    \normalsize
    \begin{itemize}
      #1
    \end{itemize}
    \vspace{0.3em}
    \textcolor{SoftInk}{课件中的论文数字均指作者在特定模型、硬件和工作负载下报告的结果，不代表无条件性能。}
  \end{frame}
}

\newcommand{\InternalCaseNote}{
  \small\textcolor{SoftInk}{说明：本页为课程研究案例或教学假设，不代表已发表论文结论；只有在实验包给出模型、硬件、工作负载和原始记录时，数字才可作为结果引用。}
}

% 14 讲细分版的统一参考资料入口。论文数字的解释边界见课程内容基线。

\newcommand{\LectureOneSources}{\ResearchSources{
  \PaperRef{vLLM: Easy, Fast, and Cheap LLM Serving with PagedAttention}{SOSP 2023}{https://doi.org/10.1145/3600006.3613165}
  \PaperRef{SGLang: Efficient Execution of Structured Language Model Programs}{arXiv:2312.07104}{https://arxiv.org/abs/2312.07104}
  \PaperRef{nano-vllm-hust 教学代码}{课程仓库}{https://github.com/intellistream/intro-to-llm-inference-engines}
}}

\newcommand{\LectureTwoSources}{\ResearchSources{
  \PaperRef{vLLM / PagedAttention}{SOSP 2023}{https://doi.org/10.1145/3600006.3613165}
  \PaperRef{Sarathi-Serve}{OSDI 2024}{https://www.usenix.org/conference/osdi24/presentation/agrawal}
  \PaperRef{DistServe}{OSDI 2024}{https://www.usenix.org/conference/osdi24/presentation/zhong-yinmin}
}}

\newcommand{\LectureThreeSources}{\ResearchSources{
  \PaperRef{vLLM / PagedAttention}{SOSP 2023}{https://doi.org/10.1145/3600006.3613165}
  \PaperRef{DistServe}{OSDI 2024}{https://www.usenix.org/conference/osdi24/presentation/zhong-yinmin}
  \PaperRef{Llumnix}{OSDI 2024}{https://www.usenix.org/conference/osdi24/presentation/sun-biao}
}}

\newcommand{\LectureFourSources}{\ResearchSources{
  \PaperRef{Orca}{OSDI 2022}{https://www.usenix.org/conference/osdi22/presentation/yu}
  \PaperRef{Sarathi-Serve}{OSDI 2024}{https://www.usenix.org/conference/osdi24/presentation/agrawal}
  \PaperRef{FastServe}{NSDI 2026}{https://www.usenix.org/conference/nsdi26/presentation/wu-bingyang}
}}

\newcommand{\LectureFiveSources}{\ResearchSources{
  \PaperRef{vLLM / PagedAttention}{SOSP 2023}{https://doi.org/10.1145/3600006.3613165}
  \PaperRef{vAttention}{arXiv:2405.04437}{https://arxiv.org/abs/2405.04437}
  \PaperRef{Rethinking KV Cache Compression}{MLSys 2025}{https://proceedings.mlsys.org/paper_files/paper/2025/hash/26289c647c6828e862e271ca3c490486-Abstract-Conference.html}
  \PaperRef{Mooncake}{FAST 2025}{https://www.usenix.org/conference/fast25/presentation/qin}
}}

\newcommand{\LectureSixSources}{\ResearchSources{
  \PaperRef{DistServe}{OSDI 2024}{https://www.usenix.org/conference/osdi24/presentation/zhong-yinmin}
  \PaperRef{RAGCache}{arXiv:2404.12457}{https://arxiv.org/abs/2404.12457}
  \PaperRef{Mooncake}{FAST 2025}{https://www.usenix.org/conference/fast25/presentation/qin}
  \PaperRef{NEO}{MLSys 2025}{https://proceedings.mlsys.org/paper_files/paper/2025/hash/66a026c0d17040889b50f0dfa650e5e0-Abstract-Conference.html}
}}

\newcommand{\LectureSevenSources}{\ResearchSources{
  \PaperRef{vLLM / PagedAttention}{SOSP 2023}{https://doi.org/10.1145/3600006.3613165}
  \PaperRef{SGLang}{arXiv:2312.07104}{https://arxiv.org/abs/2312.07104}
  \PaperRef{DistServe}{OSDI 2024}{https://www.usenix.org/conference/osdi24/presentation/zhong-yinmin}
  \PaperRef{Mooncake}{FAST 2025}{https://www.usenix.org/conference/fast25/presentation/qin}
}}

\newcommand{\LectureEightSources}{\ResearchSources{
  \PaperRef{FlashAttention}{NeurIPS 2022}{https://arxiv.org/abs/2205.14135}
  \PaperRef{FlexGen}{ICML 2023}{https://proceedings.mlr.press/v202/sheng23a.html}
  \PaperRef{PowerInfer}{SOSP 2024 / arXiv:2312.12456}{https://arxiv.org/abs/2312.12456}
  \PaperRef{FlashInfer}{MLSys 2025}{https://proceedings.mlsys.org/paper_files/paper/2025/hash/dbf02b21d77409a2db30e56866a8ab3a-Abstract-Conference.html}
}}

\newcommand{\LectureNineSources}{\ResearchSources{
  \PaperRef{DEEPSERVE}{USENIX ATC 2025}{https://www.usenix.org/conference/atc25/presentation/hu-junhao}
  \PaperRef{QServe}{MLSys 2025}{https://proceedings.mlsys.org/paper_files/paper/2025/hash/fbe2b2f74a2ece8070d8fb073717bda6-Abstract-Conference.html}
  \PaperRef{vLLM-HUST}{开源维护仓库}{https://github.com/vLLM-HUST/vllm-hust}
}}

\newcommand{\LectureTenSources}{\ResearchSources{
  \PaperRef{vLLM / PagedAttention}{SOSP 2023}{https://doi.org/10.1145/3600006.3613165}
  \PaperRef{Revisiting Disaggregated LLM Serving}{2025 preprint, arXiv:2601.08833}{https://arxiv.org/abs/2601.08833}
  \PaperRef{Rethinking KV Cache Compression}{MLSys 2025}{https://proceedings.mlsys.org/paper_files/paper/2025/hash/26289c647c6828e862e271ca3c490486-Abstract-Conference.html}
}}

\newcommand{\LectureElevenSources}{\ResearchSources{
  \PaperRef{MLPerf Inference}{MLCommons benchmark}{https://mlcommons.org/benchmarks/inference-datacenter/}
  \PaperRef{ACM Artifact Review and Badging}{ACM policy}{https://www.acm.org/publications/policies/artifact-review-and-badging-current}
  \PaperRef{Revisiting Disaggregated LLM Serving}{2025 preprint, arXiv:2601.08833}{https://arxiv.org/abs/2601.08833}
}}

\newcommand{\LectureTwelveSources}{\ResearchSources{
  \PaperRef{nano-vllm-hust}{课程教学代码}{https://github.com/intellistream/intro-to-llm-inference-engines}
  \PaperRef{vLLM-HUST}{HUST 维护的 vLLM fork}{https://github.com/vLLM-HUST/vllm-hust}
  \PaperRef{上游 vLLM}{开源项目}{https://github.com/vllm-project/vllm}
}}

\newcommand{\LectureThirteenSources}{\ResearchSources{
  \PaperRef{课程内容基线与实验材料}{课程仓库}{https://github.com/intellistream/intro-to-llm-inference-engines}
  \PaperRef{DEEPSERVE}{USENIX ATC 2025}{https://www.usenix.org/conference/atc25/presentation/hu-junhao}
  \PaperRef{GhostServe}{MLSys 2026}{https://proceedings.mlsys.org/paper_files/paper/2026/hash/d99e8e80a6c41e148db686918dd7eab3-Abstract-Conference.html}
}}

\newcommand{\LectureFourteenSources}{\ResearchSources{
  \PaperRef{课程完整材料与后续更新}{课程仓库}{https://github.com/intellistream/intro-to-llm-inference-engines}
  \PaperRef{SpecGen}{MLSys 2026}{https://proceedings.mlsys.org/paper_files/paper/2026/hash/66a026c0d17040889b50f0dfa650e5e0-Abstract-Conference.html}
  \PaperRef{GhostServe}{MLSys 2026}{https://proceedings.mlsys.org/paper_files/paper/2026/hash/d99e8e80a6c41e148db686918dd7eab3-Abstract-Conference.html}
}}



\LectureMeta{第 9 讲 异构平台与国产算力适配}{研究生课程课件}

\begin{document}

\begin{frame}[plain]
  \titlepage
\end{frame}

\begin{frame}{教学目标}
  \Core{本讲讨论平台差异如何系统性改变瓶颈位置、最优策略与验证方法。}
  \begin{itemize}
    \item 理解能运行、能优化、能验证三个层次。
    \item 理解国产异构平台的系统性差异。
    \item 讨论跨平台收益如何公平验证。
  \end{itemize}
\end{frame}

\begin{frame}{为什么平台差异是系统差异}
  \Core{不同平台在图执行、算子栈、互联和调试环境上的差异，会直接改写系统设计空间。}
  \begin{itemize}
    \item 软件栈成熟度不同。
    \item 互联协议与带宽比例不同。
    \item host-device 数据路径代价不同。
  \end{itemize}
\end{frame}

\begin{frame}{平台差异会改写哪些瓶颈}
  \Core{同一机制在不同平台上可能卡在完全不同的地方，因此不能直接平移 GPU 经验。}
  \begin{itemize}
    \item 有的平台先卡在图执行与编译稳定性。
    \item 有的平台先卡在通信与互联约束。
    \item 有的平台先卡在调试工具和观测能力不足。
  \end{itemize}
\end{frame}

\begin{frame}{适配的三个层次}
  \Core{平台适配至少分为三个层次，而真正的研究通常从第二层开始。}
  \begin{enumerate}
    \item 能运行：基本功能打通
    \item 能优化：根据平台特性重构路径
    \item 能验证：建立公平可复现的基线和对照
  \end{enumerate}
\end{frame}

\begin{frame}{从可运行到可优化的差异}
  \Core{适配工作常停留在基本可运行阶段，而研究价值主要体现在可优化与可验证两个层次。}
  \begin{itemize}
    \item 跑通只说明接口存在。
    \item 稳定收益需要重新建模瓶颈。
    \item 可验证结论需要统一实验设计与对照基线。
  \end{itemize}
\end{frame}

\begin{frame}{常见误区}
  \Core{把问题只定义成“移植工程”，往往会掩盖真正有研究价值的系统差异。}
  \begin{itemize}
    \item 只报告局部跑通，不说明收益边界。
    \item 直接沿用 GPU 经验，不重新验证假设。
    \item 用不一致 workload 比较不同平台。
  \end{itemize}
\end{frame}

\begin{frame}{跨平台实验如何更可信}
  \Core{跨平台比较最关键的是控制变量，而不是堆更多数字。}
  \begin{itemize}
    \item 统一模型、输入分布和服务目标。
    \item 显式说明软件栈、精度和并行配置。
    \item 同时报告端到端指标和关键探针指标。
  \end{itemize}
\end{frame}

\begin{frame}{跨平台比较中的常见问题}
  \Core{跨平台比较中的主要问题通常来自对照不公平、工作负载不一致或边界说明不足。}
  \begin{itemize}
    \item 输入分布不同会直接破坏比较意义。
    \item 精度和默认优化不一致会导致结论失真。
    \item 没有端到端指标时很难说明真实服务价值。
  \end{itemize}
\end{frame}

\begin{frame}{讲解：平台适配课的核心不是移植，而是重建假设}
  \Core{GPU 上形成的经验并不能直接平移到国产 NPU 或其他异构平台，因为假设条件一起变了。}
  \begin{itemize}
    \item 图执行、编译路径、通信比例、调试工具和默认优化都可能不同。
    \item 因此适配不是“接口替换”，而是重新确认瓶颈在哪、哪些机制仍然成立。
    \item 这也是为什么平台适配有研究价值，它迫使我们暴露原本被默认隐藏的系统前提。
  \end{itemize}
\end{frame}

\begin{frame}{讲解：第九讲应建立的验证边界}
  \Core{跨平台比较最容易失真，所以课程里必须先讲什么叫“可比”。}
  \begin{itemize}
    \item 模型、精度、workload、warmup 和并行配置如果不统一，数字很难解释。
    \item “能跑起来”只对应功能层，真正可研究的是“能优化”和“能验证”两层。
    \item 因此平台课的重点不是比谁更快，而是学会设计公平、可追责的比较实验。
  \end{itemize}
\end{frame}

\begin{frame}{例子：为什么同一机制在 GPU 和国产 NPU 上经常表现不同（1）}
  \Core{平台适配课的出发点不是“换了芯片”，而是执行语义、工具链和瓶颈位置一起变了。}
\begin{block}{GPU 路径}
    开发者往往更熟悉 CUDA、Nsight 和主流 kernel 栈，性能问题定位更直接。
  \end{block}
\end{frame}

\begin{frame}{例子：为什么同一机制在 GPU 和国产 NPU 上经常表现不同（2）}
  \Core{平台适配课的出发点不是“换了芯片”，而是执行语义、工具链和瓶颈位置一起变了。}
\begin{block}{NPU 路径}
    图执行、编译栈、通信路径和 profiler 语义可能完全不同，原有直觉经常失效。
  \end{block}
\end{frame}

\begin{frame}{例子：为什么同一机制在 GPU 和国产 NPU 上经常表现不同（3）}
  \Core{平台适配课的出发点不是“换了芯片”，而是执行语义、工具链和瓶颈位置一起变了。}
\begin{block}{课程结论}
    平台迁移不是简单重编译，而是一次系统假设重审。
  \end{block}
\end{frame}

\begin{frame}{例子：能跑、能优化、能验证为什么要分成三层（1）}
  \Core{很多适配工作停留在“能跑起来”，但课程必须让学生看到这和“能做研究”之间还有距离。}
\begin{block}{能跑}
    模型能返回结果，但不代表路径正确、效率合理或结果稳定。
  \end{block}
\end{frame}

\begin{frame}{例子：能跑、能优化、能验证为什么要分成三层（2）}
  \Core{很多适配工作停留在“能跑起来”，但课程必须让学生看到这和“能做研究”之间还有距离。}
\begin{block}{能优化}
    说明热点路径、编译行为和资源瓶颈已经可观察、可改动。
  \end{block}
\end{frame}

\begin{frame}{例子：能跑、能优化、能验证为什么要分成三层（3）}
  \Core{很多适配工作停留在“能跑起来”，但课程必须让学生看到这和“能做研究”之间还有距离。}
\begin{block}{能验证}
    说明你可以用公平实验解释收益，而不是把所有差异都推给平台黑盒。
  \end{block}
\end{frame}

\begin{frame}{代码例子：平台适配首先会动到哪些运行时入口（1）}
  \Core{平台适配往往不是直接改模型定义，而是先从 runner、worker、编译与内存路径切入。}
\begin{block}{典型入口}
    {\ttfamily\small
    \detokenize{model_runner.py}\par
    \detokenize{worker / executor backend}\par
    \detokenize{graph capture / compile policy}\par
    }
  \end{block}
\end{frame}

\begin{frame}{代码例子：平台适配首先会动到哪些运行时入口（2）}
  \Core{平台适配往往不是直接改模型定义，而是先从 runner、worker、编译与内存路径切入。}
\begin{block}{课堂应追问}
    这些位置分别更接近功能兼容、性能优化，还是验证与观测能力？
  \end{block}
\end{frame}

\begin{frame}{代码例子：平台适配首先会动到哪些运行时入口（3）}
  \Core{平台适配往往不是直接改模型定义，而是先从 runner、worker、编译与内存路径切入。}
\begin{block}{系统含义}
    真正的适配路径通常要先建立“哪个层面在出问题”的分层诊断。
  \end{block}
\end{frame}

\begin{frame}{代码例子：最小平台差异观察日志（1）}
  \Core{平台适配课里的代码页，重点是教学生如何建立平台差异的第一手证据。}
\begin{block}{最小日志样式}
    {\ttfamily\small
    \detokenize{log(platform, batch_size, prompt_len, stage, latency_ms)}\par
    \detokenize{log(graph_mode, compile_hit, memory_allocated)}\par
    }
  \end{block}
\end{frame}

\begin{frame}{代码例子：最小平台差异观察日志（2）}
  \Core{平台适配课里的代码页，重点是教学生如何建立平台差异的第一手证据。}
\begin{block}{课堂应追问}
    哪些量说明“路径不同”，哪些量说明“只是配置不同”？
  \end{block}
\end{frame}

\begin{frame}{代码例子：最小平台差异观察日志（3）}
  \Core{平台适配课里的代码页，重点是教学生如何建立平台差异的第一手证据。}
\begin{block}{课程结论}
    没有分层日志，平台差异就会退化成模糊叙述，无法进入系统研究语境。
  \end{block}
\end{frame}

\begin{frame}{例子：公平实验为什么在异构平台上更难（1）}
  \Core{平台比较最容易犯的错误，是把环境差异、软件成熟度差异和机制差异混在一起。}
\begin{block}{不公平一}
    比较时使用了不同精度、不同 batch 上限或不同 warmup 策略。
  \end{block}
\end{frame}

\begin{frame}{例子：公平实验为什么在异构平台上更难（2）}
  \Core{平台比较最容易犯的错误，是把环境差异、软件成熟度差异和机制差异混在一起。}
\begin{block}{不公平二}
    一方用了成熟 profiler 和手工优化，另一方仍停留在默认配置。
  \end{block}
\end{frame}

\begin{frame}{例子：公平实验为什么在异构平台上更难（3）}
  \Core{平台比较最容易犯的错误，是把环境差异、软件成熟度差异和机制差异混在一起。}
\begin{block}{课程要求}
    异构平台实验必须先交代可比边界，再谈数字输赢。
  \end{block}
\end{frame}

\begin{frame}{例子：迁移一条 decode 热路径时会遇到什么（1）}
  \Core{decode 热路径对平台差异特别敏感，因为它既要求低抖动，又持续依赖状态与调度稳定性。}
\begin{block}{算子层}
    kernel、算子融合和内存布局可能都需要重新调优。
  \end{block}
\end{frame}

\begin{frame}{例子：迁移一条 decode 热路径时会遇到什么（2）}
  \Core{decode 热路径对平台差异特别敏感，因为它既要求低抖动，又持续依赖状态与调度稳定性。}
\begin{block}{运行时层}
    graph capture、同步语义和调度节拍可能与原平台不同。
  \end{block}
\end{frame}

\begin{frame}{例子：迁移一条 decode 热路径时会遇到什么（3）}
  \Core{decode 热路径对平台差异特别敏感，因为它既要求低抖动，又持续依赖状态与调度稳定性。}
\begin{block}{研究层}
    如果不能解释是哪个层面出问题，适配工作就很难形成可发表的系统结论。
  \end{block}
\end{frame}

\begin{frame}{课堂练习：为跨平台实验写一页计划（1）}
  \Core{平台适配课最适合训练的，是写出一份短而硬的验证计划。}
\begin{block}{练习一}
    指定两个平台、一个共同 workload 和三项必须控制的变量。
  \end{block}
\end{frame}

\begin{frame}{课堂练习：为跨平台实验写一页计划（2）}
  \Core{平台适配课最适合训练的，是写出一份短而硬的验证计划。}
\begin{block}{练习二}
    指出你会优先观测的两个局部探针，以及它们为什么能帮助解释性能差异。
  \end{block}
\end{frame}

\begin{frame}{课堂练习：为跨平台实验写一页计划（3）}
  \Core{平台适配课最适合训练的，是写出一份短而硬的验证计划。}
\begin{block}{练习三}
    用一句话回答：为什么“能跑起来”还不算完成平台适配？
  \end{block}
\end{frame}


\begin{frame}{课堂讨论}
  \Core{跨平台研究需要明确平台差异究竟改写了哪些系统假设。}
  \begin{itemize}
    \item 为什么“在平台 A 上有效”不自动说明方法具有普适性？
    \item 怎样设计一个最小跨平台验证计划？
    \item 什么结论最容易因为实验不公平而失真？
  \end{itemize}
\end{frame}

\begin{frame}{本讲总结}
  \Core{平台适配本身构成系统研究的重要组成部分。}
  \begin{itemize}
    \item 平台差异会重写瓶颈与最优策略。
    \item 只有能验证，平台适配结论才有说服力。
  \end{itemize}
\end{frame}

\LectureNineSources

\begin{frame}{对应 Tutorial}
  \Core{本讲对应 Tutorial 8。先独立作答，再使用对应 Tutorial 的教师讨论要点核对。}
  \begin{itemize}
    \item 文件：\texttt{\detokenize{build/tutorials/Tutorial_08_异构平台适配.pdf}}
    \item 动手运行、日志记录与提交要求统一查看 \texttt{\detokenize{build/experiments/}} 和对应 experiment 源目录下的 \texttt{\detokenize{assignment_spring-2027/}}。
  \end{itemize}
\end{frame}

\end{document}
